\section{Decomposition into supervised agent refactors}
\label{sec:decomposition}

The roadmap of \cref{sec:roadmap} fixes phase order; this section maps
the phases onto the project's existing execution machinery. \dotnetrs{}
is developed through supervised, agent-executed refactor sessions
managed by a local orchestrator (\code{nasin-pali-pi-jan-ilo}, state in
\code{.refactor-orchestrator/}): each session binds one repository and
one branch to one written goal, a kickoff agent produces a review
document and a checklist of anchored steps
(\code{-\,[\,]\,N.M \dots\ --- refs REVIEW.md\#anchor}), each step is
executed by a fresh-context agent at a configured effort tier
(\code{draft}/\code{standard}/\code{deep}/\code{max}) and must attach
machine-checkable verification evidence, and the supervisor reviews the
per-step diff before the step commits. The repository's history
contains 51 such sessions, all merged, ranging from 6 to 68 checklist
steps. The decomposition below is stated in that system's units.

\subsection{Sizing measures}
\label{sec:decomposition-measures}

Wall-clock estimates are not meaningful for agent-executed work; the
quantities that actually predict whether a session succeeds are:

\begin{itemize}[leftmargin=1.4em]
\item \textbf{Checklist size class.} S ($\le$10 steps), M (11--25),
  L (26--45). Prior sessions above $\sim$45 steps correlate with
  plan defects (mass \emph{blocked} reports mid-session), so any
  package estimated at L or larger below is pre-split.
\item \textbf{Verification-signal quality} (\emph{high} / \emph{medium}
  / \emph{low}): whether each step has a machine-checkable acceptance
  test an agent can iterate against without supervision. Lean packages
  are the extreme high case --- a \code{sorry}-free \code{lake build}
  is a sound per-step gate, so \code{draft}-tier retry loops are
  usable for proof search. Design-type steps (choosing a
  \code{Prop}'s statement, cutting an LTS) are the low case: the
  artifact compiles regardless of whether it says the right thing,
  so tier and supervisor attention concentrate there.
\item \textbf{Step context-independence}: whether a fresh-context agent
  can execute a step from the checklist line, the review anchor, and
  \code{AGENT\_MEMORY.md} alone. Grammar-and-codegen work is highly
  independent; annotation sweeps are independent per file; model
  design is not, and such sessions need richer review documents at
  kickoff.
\item \textbf{Supervisor gate density}: the fraction of steps whose
  diff requires a substantive human judgment rather than a
  conformance glance. This is the binding resource. It is highest
  where statements (not proofs) are produced: vocabulary elaboration
  rules, axioms, LTS transitions, trust grants.
\item \textbf{Blast radius}: single-crate sessions merge cleanly;
  sessions that touch shared files across crates serialize. Prior
  experience with cascade refactors indicates slicing by crate
  dependency order, not by semantic theme.
\end{itemize}

Parallelism operates at the session level, not the step level: a
session owns one branch of one repository, so \emph{sessions in
different repositories run concurrently}, while sessions in the same
repository serialize (annotation sweeps within \dotnetrs{} could in
principle run on disjoint crates, but shared-file conflicts and the
one-branch convention argue against it). The package structure of
\cref{sec:design-compart} was chosen partly for this reason: after
WP-2, the three lanes below are independent repositories.

\subsection{Ordered work packages}
\label{sec:decomposition-wps}

Each entry is one orchestrator session (or an explicitly noted short
series), with a goal fragment of the kind the kickoff form accepts.
Lanes: \textbf{A} = \dotnetrs{}, \textbf{B} =
\path{ves-proof/crates/}, \textbf{C} =
\path{ves-proof-lean/packages/}. Within a lane, order is
strict; across lanes, an entry lists its cross-lane prerequisites.

\begin{enumerate}[leftmargin=2.2em, itemsep=0.35em]
\item \textbf{[A] Atomic-path soundness. \textbf{Done}
  (\code{208a6c8b}).} Close the unguarded
  \code{StackValue::load\_atomic} and \code{store\_atomic} paths and
  restructure the \code{SharedGlobalState} auto-trait overrides
  behind a checked discipline.
  \emph{(S; \code{deep}; high signal --- Miri leg + existing atomic
  tests.)}
\item \textbf{[A] Dynamic-gate promotion. \textbf{Done}
  (\code{208a6c8b}, same commit as WP-1).} Promote the stable Miri and
  fuzz legs from \code{continue-on-error} to blocking and extend the
  differential harness beyond its single fixture.
  \emph{(S; \code{standard}; high signal --- CI is its own oracle.)}
  Landed as \code{miri-value} and \code{fuzz-raw-memory-access} in
  \code{ci.yml}; the harness grew to seven fixtures. Phase 0 is
  therefore complete (\cref{sec:roadmap}) and lane A is unblocked up
  to WP-9, which still waits on lane B.
\item \textbf{[B/C] Repository seeds (scaffolded). \textbf{Done.}}
  Initialize the
  \code{ves-proof} Rust workspace and its five crate skeletons under
  \code{crates/}; initialize the five independent Lake package
  skeletons under \path{ves-proof-lean/packages/}; reserve the two
  manifest schemas in \path{ves-proof/schemas/}.  The base skeletons
  exist; CI and implementation work remain governed by the later
  packages.
  \emph{(S; \code{standard}; done once, largely mechanical.)}
  Verified current: both sibling repositories exist with exactly the
  reserved package layout, each source file is an inert
  \code{scaffoldOnly}/\code{SCAFFOLD\_ONLY} marker with no logic, the
  two JSON Schema drafts under \path{ves-proof/schemas/} are the only
  substantive artifacts, and no Lean toolchain is pinned yet (no
  \code{lean-toolchain}, no \code{lake-manifest.json} ---
  deliberately, per \path{ves-proof-lean/AGENTS.md}). WP-4
  (\code{ves-syntax}) is now implemented; see item 4.
\item \textbf{[B] \code{ves-syntax}. \textbf{Done}
  (\code{ves-proof} commit \code{9082659}, 2026-08-01).} Implement the
  script grammar, AST, canonical serializer, and statement/subject
  hashing against a golden corpus of scripts.
  \emph{(M; \code{standard}; high signal --- golden files and
  round-trip properties make every step self-checking.)}
  Landed the full public API: the \code{Script}/\code{Statement} AST,
  parser entry points for \code{proof!}, \code{proof\_impl!}, and
  source-text forms, the \code{canonical\_statement}/%
  \code{canonical\_subject} serializers, and the
  \code{statement\_hash}/\code{subject\_hash} SHA-256 contract with
  \code{GRAMMAR\_VERSION} pinned at \code{"0.1.0"}; 33 tests over a
  corpus of valid and malformed \code{.ves} scripts, all passing.
\item \textbf{[B] \code{ves-vocabulary}. \textbf{Done}
  (\code{ves-proof} commit \code{ad67707}, 2026-08-02).} Define the
  declarative vocabulary source for the F1--F9 constructor families and
  build the dual code generators (Rust token types, Lean declarations).
  \emph{(M; \code{deep} for the schema steps, \code{standard} for
  codegen; the elaboration-rule steps carry the highest gate density
  in lane B.)}
  Landed 71 constructors across 20 namespaces with trust-class and
  invariant-family annotations (4 trusted, 4 derived, 1 checked, 62
  static), dependency-resolution and structural
  validation (rejecting untrusted premises misclassified as trusted,
  malformed derived-dependency lists, and non-static tactics), and the
  two code generators: a \code{no\_std}-compatible Rust ghost-token
  module (\path{generated/tokens.rs}, handed to WP-6) and a Lean
  declaration inventory (\path{generated/VesVocabulary.lean}, handed to
  WP-13), both pinned by golden-file tests; \code{VOCABULARY\_VERSION}
  is now the canonical value embedded in both manifest schemas and in
  statement hashes. 17 tests, all passing. Carried forward: both
  generated artifacts are \emph{prose-lossy} with respect to the
  source. The Rust module encodes trust class structurally (a
  \code{TrustGrant} parameter, a \code{cfg} gate, dependency
  parameters) but drops the invariant family; the Lean inventory is
  \code{opaque \dots{} : Prop} declarations plus tactic-name comments,
  carrying neither trust class nor family. Consumers therefore cannot
  read vocabulary facts back out of either artifact --- see 6b.
\item \textbf{[B] \code{ves-tokens}. \textbf{Done, first pass}
  (\code{ves-proof} commit \code{73e6dea}, 2026-08-02).} Generate the
  ZST token crate:
  trust-class constructor forms, \code{cfg}-mirroring conventions, the
  \code{TrustGrant} build script, and zero-size/\code{inline} test
  matrices.
  \emph{(S--M; \code{standard}; high signal.)}
  Landed a \code{no\_std}, \code{forbid(unsafe\_code)} crate whose
  \code{lib.rs} is nothing but two \code{include!}s: the build script
  copies \path{ves-vocabulary/generated/tokens.rs} into
  \code{OUT\_DIR} and synthesizes \path{OUT_DIR/trust_grants.rs}, a
  \code{pub mod grants} holding one \code{TrustGrant} const per
  \emph{registered} trusted constructor, named in
  \code{SCREAMING\_SNAKE\_CASE} without a suffix. The registry is
  selected by the \code{VES\_TRUSTED\_TOML\_PATH} environment variable;
  an unset variable or a missing file yields an empty module, so the
  default build cannot mint any trusted premise. The test matrix is
  exhaustive by construction rather than by enumeration: a compile-time
  \code{size\_of == 0} assertion for all 72 generated types plus a test
  that re-derives the type list from the generated source and fails if
  the assertion list drifts; a test asserting every constructor is
  immediately preceded by \code{\#[inline(always)]} and that struct and
  constructor counts agree; and grant-generation tests over the golden
  vocabulary. 5 tests plus the const assertions, 6 under
  \code{--features memory-validation}; all passing.

  Three items are deliberately carried forward rather than closed, and
  each is assigned below:
  \begin{enumerate}[label=(\alph*), leftmargin=1.6em, topsep=0.2em]
  \item \emph{The \code{TRUSTED.toml} reader is a placeholder.} With no
    schema yet agreed (that is WP-9's deliverable), the build script
    registers a grant when the constructor's name occurs anywhere in
    the file as a substring, rather than parsing TOML. This is
    fail-open in the wrong direction --- a commented-out entry, or a
    name that is a substring of another, mints a grant --- so it must
    not survive into the integration bootstrap. Reassigned to WP-9,
    which now owns both halves of the contract.
  \item \emph{Trusted constructors were discovered by text-scraping.}
    The build script found them by matching
    \code{\_grant: super::TrustGrant} on the same physical line as
    \code{pub fn}, which held only because every trusted constructor
    took exactly one parameter and \code{rustfmt} kept such
    signatures on one line. A trusted-\emph{and}-derived constructor
    would have wrapped and been silently skipped. The failure was
    fail-closed (the grant is not emitted, so the citing site does not
    compile), but it was a latent trap, and the underlying cause ---
    lane B re-deriving vocabulary facts from formatted Rust --- is what
    6b removed. \textbf{Closed} by 6b below.
  \item \emph{Neither the \code{debug\_assert!} shadows nor
    lifetime-branded tokens of \cref{sec:dsl-tokens} exist.} The
    vocabulary data model has no field for a shadow predicate and no
    token is generic, so static constructors are unconditionally empty
    and F6 brands are not carried in token types. Both are v0
    simplifications, not rejections; see \cref{sec:dsl-tokens}.
  \end{enumerate}
  \smallskip
  \noindent\textbf{6b --- [B] Vocabulary manifest artifact. \textbf{Done}
  (\code{ves-proof} commit \code{c23ede8}, 2026-08-03).} Follow-up
  session, unnumbered so the cross-references above keep their
  targets. Added a third \code{ves-vocabulary} output: a canonical,
  golden-pinned, pipe-delimited manifest
  (\path{generated/vocabulary.manifest}, one row per constructor ---
  namespace, name, role, trust class, \code{cfg} feature, dependency
  list, invariant families --- under an explicit \code{manifest-format}
  version) emitted alongside the Rust and Lean artifacts. Rewired
  \code{ves-tokens}'s build script to read trusted-constructor names
  from the manifest instead of scraping \path{generated/tokens.rs},
  closing 6(b); the same file is available for WP-13 to consume so
  \code{VesCore} can distinguish a grant-axiomatized proposition from a
  derived one without re-reading the Rust.
  \emph{(S; \code{standard}; high signal --- golden files again, and
  the two consumers are existing tests.)} Landed before WP-8, as
  intended: retrofitting the trust/family distinction into an existing
  elaborator would have been more expensive than generating it up
  front. Coverage: a manifest-format guard test, a test that the
  manifest's trusted set matches the vocabulary's, and rewritten
  \code{ves-tokens} tests exercising manifest-driven discovery directly
  (rejecting an unsupported format version, rejecting a malformed row,
  an empty-registry case, and the naming-convention case), plus a new
  \code{cfg\_matrix} test asserting the manifest's recorded
  \code{cfg\_feature} for \code{checked::alignment\_guard} matches the
  crate's actual \code{memory-validation} feature gate --- the first
  machine-checked link between \cref{sec:decomposition-wps}~item~9's
  cfg-forwarding obligation and the vocabulary source. The
  \code{TRUSTED.toml} reader itself is untouched and remains WP-9's
  substring-match placeholder (\cref{sec:risks-trust}); 6b closes only
  the \emph{discovery} half of the carried-forward gap, not the
  \emph{registry-parsing} half.
\item \textbf{[B] \code{ves-macros}. \textbf{Done, first pass}
  (\code{ves-proof} commit \code{85a6abc}, 2026-08-03).} Implement
  \code{proof!}, \code{proof\_impl!}, and \code{contract!} over
  \code{ves-syntax}, with expansion-snapshot tests establishing the
  token-identity property.
  \emph{(M; \code{deep} --- proc-macro hygiene and span work; snapshot
  tests give high signal once the harness exists.)}
  Landed all four macros dispatching to the shared parser:
  expression-position \code{proof!} and item-position
  \code{proof\_impl!} expand \code{have}/\code{trust} statements to
  local \code{::ves\_tokens} bindings (the grant constant for a
  \code{trust} binding is derived by upper-casing the constructor's
  final path segment) and splice the \code{then} tokens verbatim,
  realizing token-identity by construction rather than by
  reconstruct-and-compare; \code{exempt!} validates its \code{reason}
  argument and emits nothing; \code{contract!} is an unchanged-input
  passthrough, with token generation deferred to item 10. Because
  \code{proof\_impl!} cannot introduce item-scope bindings, its
  premises type-check inside an anonymous \code{const \_: () = \{ fn
  \_\_check() \{\dots\} \};} block ahead of the verbatim \code{then}
  item --- sound, but that checker cannot capture caller-local values
  or the protected item's own generics. Three unit tests, five
  expansion-snapshot fixtures (\code{cargo-expand}/\code{macrotest}),
  and three \code{trybuild} compile-fail fixtures (a feature-gated
  checked constructor, a derived constructor with an unmet
  prerequisite, an unregistered trust path) all pass, and all three
  compile-fail cases fail as ordinary \code{rustc} type errors at the
  \code{have}/\code{trust} span rather than through
  \code{ves-macros}-specific diagnostics --- exactly the diagnostics
  model of \cref{sec:dsl-diag}. Carried forward: the crate exports bare
  \code{proof}/\code{proof\_impl}/\code{exempt}/\code{contract} macros
  rather than a \code{ves::}-namespaced facade; assembling that facade
  is integration-bootstrap work (item 9), not a \code{ves-macros}
  concern.
\item \textbf{[B] \code{ves-check}.} Implement extraction, the
  structural rules, manifest comparison, and the ratchet budgets;
  fixture-test against synthetic crates.
  \emph{(M; \code{standard}; high signal.)}
\item \textbf{[A] Integration bootstrap.} (Requires 6--8.) Add the
  \code{ves-macros}/\code{ves-tokens} dependencies, the
  \code{xtask ves-check} CI job in structural mode, and
  \code{TRUSTED.toml} seeded from the accepted-risk boundaries.
  \emph{(S; \code{standard}.)}
  Scope grew when WP-6 landed, and this is now the one session in the
  plan that must edit two repositories. It owns three things beyond
  the original list. First, the \code{TRUSTED.toml} \emph{schema} ---
  at minimum a qualified constructor name, the falsifier that monitors
  the grant (\cref{sec:risks-trust}), and an expiry or review
  condition --- together with the real parser replacing WP-6's
  substring placeholder; the parser is a lane-B edit made under a
  lane-A session, which is acceptable only because the alternative is
  a lane-B session guessing at a schema it does not consume.
  Second, the \emph{feature-propagation rule} for checked
  constructors. The central \code{cfg}-mirroring guarantee of
  \cref{sec:dsl-tokens} presumes a token constructor lives under the
  same \code{cfg} as the guard it names; the implementation instead
  gates it behind a \code{ves-tokens} Cargo feature
  (\code{memory-validation}), which is only equivalent if \dotnetrs{}
  forwards its own guard \code{cfg} to that feature rather than
  enabling it unconditionally. That forwarding must be declared in
  \dotnetrs{}'s manifest and asserted by a CI matrix leg that builds
  with the guard off and confirms the citing sites fail to compile ---
  otherwise the mechanism silently degrades to a static premise, which
  is precisely defect (a) of \cref{sec:system-defects} in a new
  costume.  Third, whether to supply the \code{debug\_assert!} shadow
  hooks of \cref{sec:dsl-tokens}: the predicates are \dotnetrs{}
  functions, so if they are wanted, the vocabulary needs a shadow field
  and \code{ves-tokens} a hook trait, and this session is where the
  requirement becomes concrete.
\item \textbf{[A] Contract layer.} Declare \code{ves::contract!} on
  the $\sim$62 documented \code{unsafe fn} and collapse the template
  comments onto contract tokens.
  \emph{(M; \code{standard}; per-function steps are fully
  context-independent.)}
\item \textbf{[A] Annotation sweep, in dependency order.} Convert
  \safetyc{} sites to scripts one crate per session:
  \crate{dotnet-utils} (M), then \crate{dotnet-value} (L,
  \emph{pre-split} into pointer/object/storage and
  stack-value/layout halves), then \crate{dotnet-runtime-memory} (M).
  Per-site steps are independent and mostly \code{standard}; sites in
  the invariant families F1/F5/F6 warrant \code{deep}.
  \emph{(4 sessions; high signal --- rustc token checking plus
  \code{ves-check} run per step.)}
\item \textbf{[C] \code{VesCore} statement layer.} (Requires 5; 6b is
  now done, so its manifest is available to tell the elaborator which
  propositions are grant-axiomatized, which are derived, and from
  what --- facts the \path{VesVocabulary.lean} inventory alone does
  not carry.)
  Implement the vocabulary mirror, the goal elaborator, statement-hash
  recomputation, and manifest ingestion, cross-checked against
  \code{ves-syntax} hashes on a shared corpus.
  \emph{(M; \code{deep} for elaboration rules --- the highest
  gate-density work in the project, since these rules are TCB;
  \code{standard} elsewhere.)}
\item \textbf{[C] $\mathcal{M}_{\mathrm{VES}}$ core.} A series of four
  sessions in fixed order: (a) memory store, alignment, and atomicity
  (\ecma{I.12.6}); (b) type layout including \ecma{II.10.7} and the
  reference-overlap rule; (c) pointer taxonomy and evaluation-stack
  typing (\ecma{II.14.4}, \ecma{III.1.1.5}, after Gordon--Syme);
  (d) reachability, finalization, and GC-handle semantics
  (\ecma{I.8.9.6.7}).
  \emph{(4 $\times$ M; \code{deep} for definitional steps,
  \code{draft}--\code{standard} for lemma backfill; medium signal
  until 14 exists --- definitions compile whether or not they are
  faithful, so supervisor review concentrates on statements.)}
\item \textbf{[C] Differential harness.} Execute
  $\mathcal{M}_{\mathrm{VES}}$ against the fixture corpus and real
  .NET; drive the model's fidelity backlog from divergences.
  \emph{(M; \code{standard}; converts lane C to high signal
  retroactively --- divergence lists become mechanical checklist
  items.)}
\item \textbf{[C] \code{VesProtocol} model.} Define the STW/arena LTS
  and model-check it exhaustively at 2--3 threads; every
  counterexample becomes a checklist item against the model or, if
  real, an issue against \dotnetrs{}.
  \emph{(M; \code{deep} for the LTS cut, \code{standard} for the
  checking harness; the LTS transitions are supervisor-heavy.)}
\item \textbf{[C] \code{RustAssumptions}.} State the
  $\mathcal{M}_{\mathrm{Rust}}$ axioms with one Miri-falsifier
  template per axiom.
  \emph{(M; \code{deep}; low proof content but every statement is
  TCB --- gate density comparable to 12.)}
\item \textbf{[C] Manifest round trip.} (Requires 8, 12.) Land the
  stub generator and proof-manifest export, and run the first
  end-to-end obligation$\to$stub$\to$proof$\to$manifest$\to$%
  \code{ves-check} cycle on a toy site.
  \emph{(S; \code{standard}; high signal --- the cycle either closes
  or does not.)}
\item \textbf{[C] Protocol theorem.} (Requires 15.) Prove the pinning
  invariants of the LTS --- the F1 keystone.
  \emph{(M--L; \code{max} with \code{draft}-tier retry loops for
  subgoals; perfect signal, hardest content; the roadmap's designated
  stall point, isolated so a stall strands nothing.)}
\item \textbf{[C] Representation relation and
  \code{GcDescFaithful}.} (Requires 13b.) Define $\mathcal{R}$ and
  prove layout faithfulness against the mechanized \ecma{II.10.7}
  rules, testing the shared corpus against the executable model first.
  \emph{(2 sessions, M each; \code{deep}; medium-to-high signal.)}
\item \textbf{[C] Site campaign.} (Requires 17; 18 for F1 sites, 19
  for F2 sites.) Retire trust grants family by family --- the two
  accepted-risk boundaries, then \code{Collect} completeness, the
  atomics family, and P/Invoke pinning --- one session per family,
  each step one lemma with the \code{sorry}-free gate.
  \emph{(4--5 $\times$ M; mixed tiers; perfect signal; steps fully
  independent, and by this point the sessions are the routine kind
  the orchestrator was built for.)}
\item \textbf{[A] Ratchet closure.} Flip \code{ves-check} from
  warning to blocking for the proved families and record the residual
  grant budget.
  \emph{(S; \code{standard}.)}
\end{enumerate}

\subsection{Pipeline shape}

The dependency structure yields three long-lived concurrent lanes with
four synchronization points: lane B is unblocked immediately after the
repository seeds (3) and runs 4$\to$8 sequentially; lane A resumes at 9 once lane
B delivers \code{ves-check} (8) --- the macros (7) already landed, so
8 is now the sole remaining gate on 9; lane C starts at 12 as soon as the vocabulary
(5) stabilizes --- which it now has (commit \code{ad67707}), so lane C
may begin while lane B continues 7$\to$8 --- and is otherwise
independent of lane A entirely; and
the only cross-lane joins are 9 (A$\leftarrow$B), 12 (C$\leftarrow$B),
17 (C$\leftarrow$B), and 20 (C$\leftarrow$A's annotation sweep, which
supplies the real obligation manifest). 6b landed inside lane B's
4$\to$8 sequence without displacing it, so join 12 now delivers the
full manifest (trust class and family included) rather than the bare
Lean inventory. Within lane C, sessions 13a--d
are sequential but 14 can begin once 13a lands, and 15/16 are mutually
independent. At peak (after 9), three sessions run concurrently ---
one per repository --- which matches the supervision model: gate
density, not agent throughput, is the binding constraint, and the
lanes are deliberately mixed so that at most one supervisor-heavy
session (12, 15, 16, 18) is active at a time while high-signal sweeps
(11, 13-backfill, 20) fill the remaining lanes.

\subsection{Current position}
\label{sec:decomposition-position}

Items 1--7 and 6b are complete: lane A's Phase 0 (1, 2), the repository
seeds (3), lane B's \code{ves-syntax}, \code{ves-vocabulary},
\code{ves-tokens}, and \code{ves-macros} (4, 5, 6, 7), and the
vocabulary-manifest follow-up (6b). The token-identity property that
makes \code{ves-syntax}'s hashes and \code{ves-tokens}'s ZSTs mean
anything at a real unsafe site is therefore realized, not just
designed, and \code{ves-tokens}'s remaining discovery mechanism ---
the \code{TRUSTED.toml} reader, not the vocabulary-manifest reader 6b
replaced --- is exactly the one placeholder explicitly reassigned to
WP-9. One session can start immediately:

\begin{itemize}[leftmargin=1.4em]
\item \textbf{Lane B: item 8, \code{ves-check}.} Now the sole gate on
  the integration bootstrap (9): lane B's sequential order (4$\to$8)
  means this is the only session lane B can run next, and it is the
  last piece \dotnetrs{} needs before it can depend on this workspace
  at all.
\item \textbf{Lane C: item 12, \code{VesCore}.} Unblocked by 5 and now
  by 6b as well, supervisor-heavy, and the longest-lead item in the
  plan. It has been runnable concurrently with lane B since revision 7
  and remains so; with 6b done, it can now read trust class and
  invariant family directly from the manifest instead of having them
  supplied by hand at elaborator-design time.
\end{itemize}

Item 8, then the integration bootstrap (9), are the only things
standing between the present state and the Phase 1 deliverable, which
is also the first point at which the study's central hypothesis ---
that scripts survive this project's refactor velocity --- is actually
under test rather than under discussion.
